\documentclass[11pt]{book}
\usepackage[utf8]{inputenc}
\usepackage{geometry}
\geometry{a5paper, margin=1.2in}
\usepackage{graphicx}
\usepackage{epigraph}
\usepackage{lettrine}
\usepackage{xcolor}
\usepackage{titlesec}
\usepackage{fontspec}
\setmainfont{Linux Libertine O}
\setsansfont{Linux Biolinum O}
\setmonofont{Inconsolata}

\titleformat{\chapter}{\huge\bfseries\scshape}{\thechapter}{20pt}{}
\titleformat{\section}{\Large\bfseries}{\thesection}{15pt}{}
\titleformat{\subsection}{\normalsize\bfseries}{\thesubsection}{12pt}{}

\newcommand{\sb}[1]{\textbf{SB #1}}
\newcommand{\dv}[1]{\textbf{DV #1}}
\newcommand{\clock}[1]{\textbf{[#1]}}

\begin{document}

\frontmatter
\title{\Huge\scshape The Reckoning Bridge\\\large A Fenwood Legacy Adventure}
\author{Found in the Fenwood ledger, the skald’s song, and a Vhasian chronicle}
\date{Year of the Everblood Rising, 876 AR}
\maketitle

\tableofcontents

\mainmatter

\chapter*{A Note on This Writing}
\addcontentsline{toc}{chapter}{A Note on This Writing}
\lettrine{T}{his} account was compiled from three voices: a Linnic raid-mother who burned the passes, a Fenwood chancellor who betrayed his king, and a Vhasian chronicler who watched his kingdom change its name. Between them, a fourth hand—perhaps the Gray Wanderer, perhaps a True Mason—has added marginal notes where the keystone rights are unclear.

The Aeler do not conquer with swords. They conquer with clauses. The bridge at Tarlington is theirs; the Tarlings only held it by sufferance. The Fenwoods learned to pay. The Tarlings did not. The road remembers.

\vspace{1em}
\noindent\textit{— The Archive of the True Masons, Tarlington chapter, autumn of 876 AR}

\clearpage

\chapter{Introduction: The Everblood Tide}

\lettrine{T}{he} Linns came from the Dolmis passes, through the old Aeler-warded gateways that the empire had once seized and the republic had neglected. They rode wolf-prowed longships up the Yloka River—the same water that the Aeler had canalized when they still held the land as a province. The bridge at Tarlington was not Tarling. It was not Viterran. It was Aeler—a keystone clause carved into every span.

The Fenwoods knew this. The Tarlings had forgotten. And the Aeler—those patient accountants of stone and breath—had been waiting for the moment to remind them.

\section*{What You Need}
\textbf{Tier:} Seasoned (41–90 XP) or Veteran (91–150 XP) \quad \textbf{Players:} 3–6 \quad \textbf{Length:} 3–4 sessions

\subsection*{Core Clocks}
\begin{itemize}
    \item \textbf{Linnic Momentum [8]} – Advances as the invasion succeeds. When full, the Linns reach Tarlington before the Fenwoods.
    \item \textbf{Tarling Resistance [6]} – Advances when the Tarlings rally. When full, the invasion stalls; the Fenwoods must negotiate or be crushed between armies.
    \item \textbf{Fenwood Reckoning [6]} – The countdown to the Fenwoods’ own betrayal. Advances when they give aid to the Linns or seize Aeler infrastructure. When full, they must declare their allegiance openly.
    \item \textbf{Keystone Debt [6]} – Advances when the party takes Aeler aid, fails to offer tribute, or damages infrastructure. When full, an Aeler factor declares the bridge closed to all Fenwood traffic.
    \item \textbf{True Mason Favor [4]} – Advances when the party aids a repair effort, respects Aeler rites, or offers an heirloom. When full, the Masons reinforce the bridge; it cannot be burned this session.
\end{itemize}

\subsection*{Aeler Factions}
\begin{itemize}
    \item \textbf{True Masons} – Mendicant order of masonry mages. They repair and maintain Aeler infrastructure. They respect tradition and payment.
    \item \textbf{Fortune-Delvers} – Rogue prospectors who seek profit. They are not bound by Masonic oaths and will side with the highest bidder.
    \item \textbf{Aeler Factors} – Guild officials who collect tolls, audit infrastructure, and hold keystone rights. They are distant; they deal in ledgers, not loyalty.
\end{itemize}

\subsection*{Tone}
This is not a war of glory. It is a war of ledgers, keystones, and unpaid debts. The Linns bring fire; the Aeler bring contracts; the Fenwoods bring a hundred years of humiliation sharpened into steel. Play not for honor, but for leverage. And remember: \textit{the road does not forgive a debt. It only reschedules it.}

\clearpage

\chapter{Act I: The Corsair's Map}
\textbf{Narrator: Captain S?ren the One-Eyed, Linnic Raid-Mother}

\lettrine{W}{e} knew the crossing before we saw it. The Aeler built a bridge of stone and shadow across the Tarlington Gap—a span so broad that market stalls line its middle, so high that the river below is a whisper. The old emperors claimed it as their own, but the keystone still bears the mark of the Aeler Factor. They never forget a debt.

“The bridge is neutral ground,” a True Mason told me, years ago. “The Tarlings hold it by Aeler grace. That grace can be revoked.”

We came down the Yloka at the spring thaw. The ice broke behind us like a cheering army. Ahead, the bridge waited.

\section*{Opening Scene: The Reed Camps}
The PCs are Linnic raiders (or Fenwood allies) encamped at the headwaters of the Yloka. Longships are hidden in reed beds. Scouts have mapped every ford.

Read aloud:
\begin{quote}
\emph{The ice breaks on the Yloka. S?ren the One-Eyed spreads a chart across a drum—skin from a Utaran tax collector, still bearing the imperial grid. “Tarlington is the hinge,” she says. “Take the bridge, and the river roads are ours. The Tarlings will burn their own fields to stop us.” She looks at you—a Fenwood agent, perhaps—and smiles. “Your father’s iron kept our keels afloat. Now your steel must keep our feet dry.”}
\end{quote}

\subsection*{Clock: Linnic Momentum [8]}
Advances by 1 when the party secures a ford, defeats a Tarling patrol, or delivers supplies. When full, the Linns reach Tarlington first.

\subsection*{Key Encounter: The True Mason by the Road}
As the party scouts the northern approach, a lone figure in grey robes sits on a milestone, tapping a stone rod against the pavement.

Read:
\begin{quote}
\emph{The man does not rise. “My name is Orin of the True Masons. I am not a soldier. I am not a diplomat. I repair things.” He taps the rod again. The stone hums. “This bridge was consecrated to the Aeler Compact. No king, no jarl, no chancellor may close it without paying the toll. So I ask: who will pay?”}
\end{quote}

Orin offers to keep the bridge neutral—for a price. The party must:
\begin{enumerate}
    \item Provide a \textbf{Keystone Offering} – a valuable item (heirloom, magic, or promise) to be held by the Aeler guild.
    \item \textbf{Swear the Bridge Oath} – All parties (Fenwood, Linns, Tarlings) agree to arbitration by an Aeler factor. (This prevents the bridge from being burned or permanently seized.)
\end{enumerate}
If the party agrees, the bridge becomes a \textbf{Safe Zone} – no fighting on the span. The battle shifts to the shorelines. Gain \textbf{True Mason Favor +2}. If they refuse, Orin sighs and vanishes. The bridge’s wards weaken; the Tarlings may attempt to burn it (add \textbf{Bridge Fire [4]} clock).

\subsection*{Key Encounter: The Fortune-Delver}
Later, the party encounters an Aeler \textbf{Fortune-Delver}—a prospector of opportunity, not bound by the Masons’ oaths.

Read:
\begin{quote}
\emph{She appears in the Fenwood camp, uninvited, stirring a pot of thin stew. “I hear you’re racing the Linns to Tarlington. I can show you a route the Aeler engineers used—a maintenance shaft that bypasses the gates. But the keystone seal is intact. You’ll need my key.” She holds up a brass rod. “Cost: a share of the bridge’s tolls for the next ten years.”}
\end{quote}

\textbf{Offer:} Gain \textbf{+2} to \textbf{Fenwood Leverage} (see Act II) if the party accepts. Cost: a \textbf{Debt Clock [6]} to the Aeler syndicate. If they refuse, she shrugs and sells the same information to the Linns (adds \textbf{Linnic Momentum +2}).

\subsection*{Act I Resolution}
The party reaches the Viterran plain. The Tarlings have fortified the eastern roads. The Fenwood agents must choose where to send the Linnic scouts.

\textbf{Decision Point:} Reveal the secret supply routes (gains \textbf{Linnic Momentum +2}, \textbf{Fenwood Reckoning +2}) or hold back (gain \textbf{Tarling Resistance –1} but risk the Linns’ trust).

\clearpage

\chapter{Act II: The Chancellor's Gambit}
\textbf{Narrator: Aldric Fenwood, Chancellor of the Fens}

\lettrine{M}{y} father took the name Fenwood so we might live. I have spent thirty years building granaries, forging contracts, smiling at Tarling feasts while my brothers sharpened blades in the fens. The Linns did not find our iron by accident. I sent it to them—first in trickles, then in wagon trains disguised as lumber.

The Tarlings do not know. They call me \textit{Chancellor}, seat me at the high table, speak of my loyalty to the crown. They do not know that the crown is already cracked.

The Aeler know. They watch from their toll-keeps. A Fortune-Delver once told me: “The bridge is a ledger. Every crossing is an entry. You cannot erase a line; you can only pay it forward.”

I intend to pay.

\section*{The Fenwood Estate}
The party arrives at the Fenwood keep (now improved from the pithy county of the Exile Road). It is a grim place of black timber and whitewashed stone, surrounded by drained fields and peat stores.

Read:
\begin{quote}
\emph{The Fenwood hall is cold despite the hearth. Aldric—grey-haired, with the scars of a Black Banner youth—does not rise. “The Linns have broken the eastern passes,” he says. “Tarlington will fall within the month. I have a choice to make. So do you.”}
\end{quote}

\subsection*{Clock: Tarling Suspicion [4]}
Advances when the PCs are seen with Linnic agents, or when Aldric’s secret is revealed. When full, the Tarlings arrest Aldric and seize the keep. \textbf{Fenwood Reckoning} resets to 0, and the party must rescue him.

\subsection*{Key Encounter: The Tarling Council}
The Tarling duke summons the Fenwoods to a war council. He demands a levy of men and iron. Aldric must decide how much to give.

Read:
\begin{quote}
\emph{The Tarling duke grips his throne. “The heathens are at our gates. Every lord owes me swords. What say you, Fenwood?” Behind him, a map of Viterra shows half the province already in flames.}
\end{quote}

\textbf{Options:}
\begin{itemize}
    \item \textbf{Loyalty Feigned} – Send a token force. \textbf{Tarling Suspicion +1}, \textbf{Linnic Momentum –1}.
    \item \textbf{Honest Aid} – Commit real troops. \textbf{Tarling Suspicion –1}, \textbf{Fenwood Reckoning +2}.
    \item \textbf{The False Levy} – Send men secretly loyal to the Fenwoods, dressed as Tarling conscripts. \textbf{Wits + Deception}, DV 4. On success, the troops will turn at a signal. \textbf{Fenwood Reckoning +1}.
\end{itemize}

\subsection*{Key Encounter: The True Mason’s Audit}
A True Mason arrives at the keep to inspect the keystone rights of the Fenwood granaries. (The Aeler hold the original title to the land.)

Read:
\begin{quote}
\emph{A woman in a dust-grey cloak unrolls a parchment. “The Compact requires an annual tithe in grain or silver. Your family has not paid in three years.” She taps the parchment. “The debt is now twelve hundred bushels. Or you may pledge a service.”}
\end{quote}

\textbf{Options:}
\begin{itemize}
    \item Pay the grain. (\textbf{Survival –3}, but \textbf{True Mason Favor +2})
    \item Pledge a service: repair the Tarlington Bridge after the war. (Gain \textbf{Obligation [4]} to the Masons, but \textbf{Keystone Debt –2})
    \item Refuse. (\textbf{Keystone Debt +2}; the Mason warns that the bridge may be closed to Fenwood traffic.)
\end{itemize}

\subsection*{The Linnic Envoy}
A Linnic captain (S?ren’s kin) arrives at the keep under cover of night. She demands to know when the Fenwoods will strike. Aldric reveals his plan: the Fenwoods will march on Tarlington before the Linns reach it, seizing the bridge and claiming the conquest for themselves—then negotiate from strength.

Read:
\begin{quote}
\emph{“You would rob us of the prize?” the captain snarls. “We bled for the passes. Our dead freeze in the hills.”}
\emph{Aldric pours mead. “You seek to be kings of Viterra. I seek to be the one who opens the gate. We are not enemies. We are partners. The bridge is large enough for both our banners.”}
\end{quote}

The party must decide: support Aldric’s plan (gain \textbf{Fenwood Reckoning +3}, risk alienating the Linns) or cede the glory to the Everbloods (gain \textbf{Linnic Momentum +3}, but the Fenwoods become servants, not allies).

\subsection*{Act II Resolution}
Record \textbf{Fenwood Leverage} as a new clock (starting 0). It advances when the party secures Aeler aid, seizes keystone rights, or outmaneuvers the Linns. When full (4), the Fenwoods control the bridge’s center.

\clearpage

\chapter{Act III: The Bridge of Stone and Shadow}
\textbf{Narrator: Brother Mavron, Vhasian Chronicler}

\lettrine{I}{} was in Tarlington when the news came. The Linns had crossed the Yloka. The Fenwoods had marched from the fens. And the great bridge—the Spine of Viterra, which had stood since the Aeler first laid its keystone—became a battlefield of three banners.

The Tarlings burned the western span to slow the invaders. The fire caught the old wards. For one night, the bridge existed in two places at once—as stone and as shadow. I saw men walk into the mist and emerge hours older, or not at all.

The Aeler wraiths rose from the river. They were not ghosts of the empire. They were older—echoes of the True Masons who had first laid the keystone. They did not mourn the Tarlings. They did not fear the Linns. They enforced the Compact.

I wrote what I saw. I do not claim to understand it.

\section*{The Three-Way Battle}
The Tarlings hold the eastern gatehouse. The Linns press from the north. The Fenwoods approach from the south, their true colors still hidden. Each faction has a representative (PC or NPC) who can negotiate or fight.

\subsection*{Clocks}
\begin{itemize}
    \item \textbf{Tarling Resolve [6]} – Advances with Tarling successes. When full, they refuse all terms and fight to the death.
    \item \textbf{Linnic Fury [6]} – Advances with Linnic brutality. When full, they storm the bridge regardless of negotiation.
    \item \textbf{Fenwood Leverage [4]} – Advances when the PCs reveal secrets, sabotage, or rally allies. When full, they control the bridge’s center.
\end{itemize}

\subsection*{The Wraiths’ Price}
The ancient bridge-wardens appear, translucent figures in chiseled armor. They speak in unison:

Read:
\begin{quote}
\emph{“The Spine wakes. Name your crossing or turn back. The living do not dictate terms to the dead. The Compact is eternal. Who will pay the toll?”}
\end{quote}

The wraiths accept payment in concrete terms:
\begin{itemize}
    \item \textbf{The Laurel Seal} (if still held) – Surrender to the Aeler guild. Gain \textbf{True Mason Favor +4}, but lose the heirloom permanently.
    \item \textbf{A Repair Vow} – Swear to rebuild a collapsed Aeler aqueduct in the fens. \textbf{Obligation [6]}, but the bridge stands.
    \item \textbf{A Keystone Offering} – A magical item or a significant treasure. The wraiths test its value by weighing it against the bridge’s shadow. \textbf{Wits + Lore}, DV 4.
\end{itemize}
If the party fails to pay, the wraiths withdraw their protection. The Tarlings may now burn the span (add \textbf{Bridge Fire [4]} clock).

\subsection*{The Fortune-Delver’s Double-Cross}
If the party accepted her aid in Act I, she now appears on the bridge’s midpoint, flanked by Aeler factors.

Read:
\begin{quote}
\emph{“The Compact has been updated,” she announces. “The bridge will remain open—for a toll. One silver per head. And the Fenwoods will collect it on our behalf. That was the deal.” She smiles. “You didn’t ask whose share I was taking.”}
\end{quote}

\textbf{Outcome:} The party gains control of the toll (a permanent \textbf{Standard Asset} for Fenwood income), but the Aeler take a 40% cut (reduce profit). They may attempt to renegotiate (\textbf{Presence + Sway}, DV 5) or refuse and fight (the factors are guarded by \textbf{Cap 3} Aeler armsmen).

\subsection*{The Betrayal}
At the critical moment, Aldric Fenwood signals his hidden troops. The Fenwood levy turns on the Tarling rear guard. The gatehouse falls.

Read:
\begin{quote}
\emph{The Tarling duke stares as his own knights lower their lances at his back. “Fenwood,” he whispers. “I gave you land. I gave you bread.”}
\emph{Aldric does not look away. “You gave us humiliation. The road does not forgive. It reschedules.”}
\end{quote}

\textbf{Choice for the PCs:}
\begin{itemize}
    \item \textbf{Mercy} – Spare the duke. Gain \textbf{Obligation [4]} to a surviving Tarling heir. \textbf{Fenwood Reckoning –2}.
    \item \textbf{Execution} – Kill him. Future \textbf{Tarling Favor –4}, \textbf{Linnic Momentum +2}.
\end{itemize}

\clearpage

\chapter{Act IV: The King’s Capitulation}
\textbf{Narrator: Chancellor Aldric Fenwood (final confession)}

\lettrine{T}{he} war ended not on the bridge, but in a tent on the south bank. The King of Vhasia—a Blood, a man whose ancestors had wed Valvano princesses—came to parley. He offered the Linns what they wanted: the crown of Viterra, the title of Duke in Vhasia, and a marriage pact to seal it.

They called it a capitulation. It was a coronation. The Everbloods rose that day, and the Tarlings sank into the marshes of memory.

As for me, I knelt to the new king. He called me “Chancellor.” I called him “my lord.” We both knew the debt was unpaid. The road only reschedules.

\section*{The Feast of Two Thrones}
The Linnic captain (now Jarl Everblood) takes the crown. The King of Vhasia signs the treaty. The Fenwoods are granted the title of Count in their own lands and a place at the high table.

\subsection*{Outcomes}
\begin{itemize}
    \item \textbf{If the Fenwoods seized the bridge} – They are named \textbf{Marquesses of the Tarlington Gap}, with control of the Aeler bridge’s southern tolls. \textbf{Fenwood Leverage} becomes a permanent \textbf{Standard Asset}.
    \item \textbf{If they ceded to the Linns} – They are named \textbf{Stewards of the Fens}, with no military power but a marriage alliance to the new dynasty. Gain \textbf{Everblood Promise} (String).
    \item \textbf{If the Aeler factor was honored} – The True Masons grant the Fenwoods a \textbf{Keystone Charter}. Once per session, they may reroute a toll or bridge crossing, treating \textbf{Desperate} as \textbf{Controlled}.
\end{itemize}

\subsection*{The Ledger Closes}
Aldric returns to his keep. He writes in the leather book—the same ledger his ancestors carried from Ecktoria.

Read:
\begin{quote}
\emph{“We buried the Valvano name so we could live. Today, we helped bury the Tarlings so we could rise. The road is not straight. It is not kind. It is only long.}
\emph{Let the fourth sack come. We will be ready.”}
\end{quote}

\clearpage

\chapter*{Epilogue: The Chancellor’s Ledger}
\addcontentsline{toc}{chapter}{Epilogue: The Chancellor's Ledger}

\lettrine{A}{} True Mason arrives at the Fenwood keep a month after the battle. He inspects the keystone of the new granary, nods, and stamps a slate tile.

“The Compact holds,” he says. “The bridge will stand for another thousand years—if you keep your promises.”

Aldric signs his name in the Aeler ledger—not as Valvano, not as a Tarling vassal, but as \textit{Chancellor of the Fens and Warden of the Spine}. The road reschedules.

\vspace{2em}
\noindent\textit{— From the Tarlington archives, copied in the year 876 AR, under the lamp, bread and salt. The ink is ash and river-mud.}

\clearpage

\chapter*{Appendix: Followers & Assets}
\addcontentsline{toc}{chapter}{Appendix: Followers & Assets}

\section*{True Mason Apprentice}
If the party aided the True Masons, they may gain a young Mason as a follower.

\textbf{Cap: 2} \quad \textbf{Upkeep: 1} (food and materials)
\begin{itemize}
    \item \textbf{Independent Action (once per session):} \textit{“He reads the stone.”} – Diagnose structural weakness or hidden passages. Gain \textbf{+2 dice} to any roll involving infrastructure or collapse.
    \item \textbf{Harm: 2} \quad \textbf{Exposure: 3}
\end{itemize}

\section*{The Bridge-Wraith’s Boon (Epic Asset, 12 XP)}
If the PCs appeased the Aeler wraiths, one of them may gain a fragment of their oath.
\begin{itemize}
    \item \textbf{Free Use (once per session):} Declare that a doorway or passage is “bound.” Creatures with supernatural origins cannot cross it for one exchange.
    \item \textbf{Boon Activation (1 Boon):} Ask the wraith a question about a bridge or crossing. It answers truthfully.
\end{itemize}

\section*{Everblood Marriage Pact (Standard Asset)}
If the Fenwoods married into the new dynasty.
\begin{itemize}
    \item \textbf{Free Use (once per session):} Gain \textbf{+1 die} on social rolls with any Everblood-affiliated court.
    \item \textbf{Boon Activation (1 Boon):} Call in a minor favor from a distant cousin—supplies, safe passage, or a distraction.
\end{itemize}

\section*{Keystone Charter (Standard Asset)}
If the Fenwoods honored the Aeler Compact.
\begin{itemize}
    \item \textbf{Free Use (once per session):} Reroute a toll or bridge crossing. Treat a \textbf{Desperate} navigation roll as \textbf{Controlled}.
    \item \textbf{Boon Activation (1 Boon):} Temporarily close a bridge to a pursuing enemy for one exchange.
\end{itemize}

\clearpage

\chapter*{Quick Reference: Clocks}
\addcontentsline{toc}{chapter}{Quick Reference: Clocks}

\begin{tabular}{|l|l|l|l|}
\hline
Clock & Size & Advances When & Filled Effect \\
\hline
Linnic Momentum & 8 & Securing fords, defeating patrols, delivering supplies & Linns reach Tarlington first \\
\hline
Tarling Resistance & 6 & Tarlings rally, successful counter-attacks & Invasion stalls; Fenwoods caught between \\
\hline
Fenwood Reckoning & 6 & Aiding Linns, seizing bridge infrastructure & Must declare allegiance openly \\
\hline
Keystone Debt & 6 & Taking Aeler aid, refusing tribute, damaging wards & Aeler factor closes bridge to Fenwoods \\
\hline
True Mason Favor & 4 & Offering heirlooms, aiding repairs, respecting rites & Masons reinforce bridge; cannot be burned \\
\hline
Tarling Suspicion & 4 & Seen with Linns, secrets exposed & Aldric arrested; rescue mission required \\
\hline
Fenwood Leverage & 4 & Securing Aeler aid, outmaneuvering rivals & Control of bridge’s center \\
\hline
Bridge Fire & 4 & Wards weakened, Tarlings desperate & Span burns; battle shifts to the river \\
\hline
\end{tabular}

\clearpage
\chapter*{Colophon}
\addcontentsline{toc}{chapter}{Colophon}

\lettrine{T}{his} book was written in a toll-keep on the south bank of the Tarlington Gap, on paper that had once been a customs manifest. The ink is lamp-black cut with the ash of a burned Aeler contract. The binding is stone-grey linen.

The Gray Wanderer notes: the Aeler are “soft colonizers”—they do not wave banners, but their keystones are everywhere. The Fenwoods learned to pay. The Tarlings did not. The road remembers.

\vspace{1em}
\noindent The ninth cup remains un-poured. The ninth name remains unspoken.

\end{document}
